\begin{table}[!htbp]
\centering\footnotesize
\caption{Subamostragem no treinamento: análise exploratória, posterior à avaliação principal. Ordenação isotônica abaixo do corte.}
\label{tab:cap8-subamostragem}
\setlength{\tabcolsep}{3pt}
\begin{tabular}{>{\raggedright\arraybackslash}p{0.19\linewidth}>{\raggedleft\arraybackslash}p{0.15\linewidth}>{\raggedleft\arraybackslash}p{0.11\linewidth}>{\raggedleft\arraybackslash}p{0.11\linewidth}>{\raggedleft\arraybackslash}p{0.16\linewidth}>{\raggedleft\arraybackslash}p{0.16\linewidth}}
\toprule
Pos.:neg. & \shortstack[r]{PAR=2 no\\topo 96} & Brier & ECE & \shortstack[r]{Sens. topo 96\\passo 11 (\%)} & \shortstack[r]{Sens. topo 96\\sexo M (\%)} \\
\midrule
Sem subamostragem & 22 & 0{,}000629 & 0{,}000588 & 64{,}0 & 57{,}7 \\
1:1 & 13{,}0 [8; 17] & 0{,}000999 & 0{,}001161 & 38{,}0 & 25{,}0 \\
1:5 & 18{,}5 [16; 21] & 0{,}000806 & 0{,}000811 & 54{,}0 & 46{,}2 \\
1:20 & 21{,}0 [18; 22] & 0{,}000698 & 0{,}000685 & 64{,}0 & 53{,}8 \\
1:50 & 22{,}5 [20; 24] & 0{,}000648 & 0{,}000623 & 66{,}0 & 59{,}6 \\
\bottomrule
\end{tabular}
\par\smallskip\parbox{\linewidth}{\footnotesize Medianas entre sementes; colchetes indicam mínimo e máximo, não IC. A linha Sem subamostragem é o treino sem subamostragem, executado uma vez. A avaliação mantém a prevalência original. Sens. topo 96: sensibilidade no topo 96 da ordenação, no grupo indicado. Métricas das outras variantes e grupos: \path{data/subamostragem_treino_resumo.csv}.}
\par\footnotesize A ordenação usa a probabilidade calibrada do classificador. A linha Sem subamostragem corresponde ao treino sem subamostragem, não à referência operacional de pares verdadeiros.
\end{table}
